\documentclass[8pt]{beamer}

\usepackage[english]{minimalistslides}

\setthemename{Southeast University}
\setthemelogo{figs/logo.jpg}

% Project metadata: most users only need to update this block.
\renewcommand{\presentationtitle}{Minimalist Academic Slides}
\renewcommand{\presentationauthor}{Zhang San}
\renewcommand{\presentationadvisor}{Li Si}
\renewcommand{\presentationinstitution}{\themename}
\renewcommand{\presentationdate}{\today}

\begin{document}

% Two cover styles are included so users can keep the one they prefer.
\makeboxedcover
\makelinecover

\begin{frame}{Outline}
    \tableofcontents
\end{frame}

\section{Getting Started}

\makesection

\subsection{What Is Included}

\begin{frame}[fragile]{\insertsection}{\insertsubsection}

    This repository packages a clean \LaTeX{} Beamer workflow for academic presentations. The default theme is inspired by Southeast University, but the structure is intentionally lightweight so it can be rebranded quickly.

    \begin{table}
        \small
        \renewcommand{\arraystretch}{1.4}
        \centering
        \begin{tabular}{ll}
            \toprule
            File Path & Purpose \\
            \midrule
            \lstinline|main-en.tex| & English demo deck and authoring entry point \\
            \lstinline|main-zh.tex| & Chinese demo deck on the same branch \\
            \lstinline|minimalistslides.sty| & Shared package for theme, layout, and commands \\
            \lstinline|reference.bib| & Bibliography entries used by the sample slides \\
            \lstinline|figs/logo.jpg| & Institution logo shown in the frame header \\
            \lstinline|figs/example.jpg| & Example image asset used in the demo \\
            \bottomrule
        \end{tabular}
    \end{table}

\end{frame}

\subsection{Retheme in Minutes}

\begin{frame}[fragile]{\insertsection}{\insertsubsection}

    \begin{block}{Fast customization path}
        \begin{enumerate}
            \item Replace \lstinline|figs/logo.jpg| with your organization logo.
            \item Update the metadata block near the top of \lstinline|main-en.tex|.
            \item Adjust the palette with \lstinline|\setthemeprimaryrgb| and \lstinline|\setthemeaccentrgb| if needed.
            \item Reuse the same style package in your own deck instead of editing the package internals first.
        \end{enumerate}
    \end{block}

    \begin{callout}{Why this structure works}
        Theme behavior now lives in a single package, while language-specific content lives in separate entry files. That keeps customization easy without duplicating style logic.
    \end{callout}

\end{frame}

\subsection{Build Workflow}

\begin{frame}[fragile]{\insertsection}{\insertsubsection}

    \begin{block}{Recommended build command}
        The repository includes a \lstinline|.latexmkrc| file, so a single command is enough for the full build loop.
        \begin{lstlisting}[style=latex]
latexmk -xelatex main-en.tex
        \end{lstlisting}
    \end{block}

    \begin{block}{Manual compilation}
        If you prefer explicit steps, use the following sequence.
        \begin{lstlisting}[style=latex]
xelatex main-en.tex
bibtex main-en
xelatex main-en.tex
xelatex main-en.tex
        \end{lstlisting}
    \end{block}

\end{frame}

\section{Authoring Slides}

\makesection

\subsection{Sections and Outline Slides}

\begin{frame}[fragile]{\insertsection}{\insertsubsection}

    \begin{block}{Sections}
        \begin{enumerate}
            \item Use \lstinline|\section{}| and \lstinline|\subsection{}| to organize your deck.
            \item Use \lstinline|\makesection| after each new section if you want a divider slide.
            \item Adjust the divider width with \lstinline|\makesection[0.5]| when needed.
        \end{enumerate}
    \end{block}

    \begin{lstlisting}[style=latex]
\section{Method}
\subsection{Training Pipeline}
\makesection
    \end{lstlisting}

    \begin{block}{Outline slide}
        Use a regular frame with \lstinline|\tableofcontents| to generate the slide outline.
    \end{block}

\end{frame}

\subsection{Frames and Reusable Blocks}

\begin{frame}[fragile]{\insertsection}{\insertsubsection}

    \begin{block}{Frame types}
        \begin{enumerate}
            \item \lstinline|\begin{frame}{title}{subtitle}| is the default pattern.
            \item \lstinline|\begin{frame}[plain]| is useful for cover or closing slides.
            \item \lstinline|\begin{frame}[fragile]| is required when a slide contains code blocks.
        \end{enumerate}
    \end{block}

    \begin{block}{Custom block}
        A standard content block looks like this.
        \begin{lstlisting}[style=latex]
\begin{block}{Key Insight}
    Put your main argument here.
\end{block}
        \end{lstlisting}
    \end{block}

\end{frame}

\begin{frame}[fragile]{\insertsection}{\insertsubsection}

    \begin{block}{Accent blocks}
        You can temporarily switch the semantic color to create contrast for warnings, emphasis, or side notes.
    \end{block}

    \colorlet{themecolor}{themered}
    \begin{block}{Accent Example}
        This block uses the accent color without changing the rest of the deck.
        \begin{lstlisting}[style=latex]
\colorlet{themecolor}{themered}
\begin{block}{Accent Example}
    ...
\end{block}
\colorlet{themecolor}{themegreen}
        \end{lstlisting}
    \end{block}
    \colorlet{themecolor}{themegreen}

\end{frame}

\subsection{Lists and Text Emphasis}

\begin{frame}{\insertsection}{\insertsubsection}

    \begin{block}{Lists}
        \begin{itemize}
            \item The theme keeps bullets minimal and color-coded.
            \item Nested lists remain readable for dense academic material.
            \begin{itemize}
                \item First nested level
                \begin{itemize}
                    \item Second nested level
                \end{itemize}
            \end{itemize}
        \end{itemize}

        \begin{enumerate}
            \item Ordered lists are useful for procedures and experiments.
            \item Nested enumerations inherit lighter visual weight.
        \end{enumerate}
    \end{block}

\end{frame}

\begin{frame}[fragile]{\insertsection}{\insertsubsection}

    Footnotes are supported as usual with \texttt{\textbackslash footnote\{\}}\footnote{This is a sample footnote.}, and the table below summarizes a few useful emphasis helpers exposed by the theme.

    \begin{table}
        \small
        \renewcommand{\arraystretch}{1.4}
        \setlength{\tabcolsep}{2mm}
        \begin{tabular}{llll}
            \toprule
            Command & Output & Command & Output \\
            \midrule
            \lstinline|\textsf{}| & \textsf{Sans-serif} & \lstinline|\textbf{}| & \textbf{Bold} \\
            \lstinline|\textrm{}| & \textrm{Serif} & \lstinline|\texttt{}| & \texttt{Monospace} \\
            \lstinline|\uline{}| & \uline{Underline} & \lstinline|\sout{}| & \sout{Strikethrough} \\
            \lstinline|\highlight{}| & \highlight{Highlight} & \lstinline|\stronghighlight{}| & \stronghighlight{Strong highlight} \\
            \bottomrule
        \end{tabular}
    \end{table}

    \begin{myquote}
        This quotation environment is intentionally understated, so it also works well for notes, caveats, and reviewer feedback.
    \end{myquote}

\end{frame}

\section{Content Examples}

\makesection

\subsection{Figures and Tables}

\begin{frame}[fragile]{\insertsection}{\insertsubsection}

    \begin{block}{Figures}
        Standard \LaTeX{} figures work without any custom wrapper.
        \begin{lstlisting}[style=latex]
\begin{figure}
    \centering
    \includegraphics[width=0.55\textwidth]{figs/example.jpg}
    \caption{Example figure}
\end{figure}
        \end{lstlisting}
    \end{block}

    \begin{figure}
        \centering
        \includegraphics[width=0.45\textwidth]{figs/example.jpg}
        \caption{Example figure}
        \label{fig:example}
    \end{figure}

\end{frame}

\begin{frame}[fragile]{\insertsection}{\insertsubsection}

    \begin{block}{Tables}
        \begin{lstlisting}[style=latex]
\begin{table}
    \centering
    \caption{Sample table}
    \begin{tabular}{cccc}
        \toprule
        Col 1 & Col 2 & Col 3 & Col 4 \\
        \midrule
        ...
    \end{tabular}
\end{table}
        \end{lstlisting}
    \end{block}

    \begin{table}
        \centering
        \caption{Sample table}
        \label{tab:example}
        \begin{tabular}{cccc}
            \toprule
            Column 1 & Column 2 & Column 3 & Column 4 \\
            \midrule
            1 & 2 & 3 & 4 \\
            5 & 6 & 7 & 8 \\
            \bottomrule
        \end{tabular}
    \end{table}

\end{frame}

\subsection{Math and Code}

\begin{frame}[fragile]{\insertsection}{\insertsubsection}
    \label{frame:math-and-code}

    \begin{block}{Equations}
        The theme stays out of the way so mathematical notation remains standard.
        \begin{equation}
            \label{eq:example}
            i\hbar\frac{\partial \psi}{\partial t}
            = \frac{-\hbar^2}{2m} \left(
            \frac{\partial^2}{\partial x^2}
            + \frac{\partial^2}{\partial y^2}
            + \frac{\partial^2}{\partial z^2}
            \right) \psi + V\psi
        \end{equation}
    \end{block}

    \begin{block}{Code listings}
        Use \lstinline|\lstlisting| for larger snippets and \lstinline|\lstinline| for inline examples.
        \begin{lstlisting}[style=python]
def build_slides(title, author):
    return {
        "title": title,
        "author": author,
        "engine": "xelatex",
    }
        \end{lstlisting}
    \end{block}

\end{frame}

\subsection{References and Cross-References}

\begin{frame}[fragile]{\insertsection}{\insertsubsection}

    \begin{block}{Bibliography}
        Keep references in \lstinline|reference.bib| and cite them inline when needed. For example, Kopka et al. remain a classic introduction to \LaTeX{}\footfullcite{kopka2003guide}.
    \end{block}

    \begin{block}{Cross-references}
        Figures, tables, equations, and frames can be linked with the standard \lstinline|\ref{}| and \lstinline|\pageref{}| commands. For example, equation (\ref{eq:example}) appears on page \pageref{frame:math-and-code}, and Figure \ref{fig:example} shows the sample image asset.
    \end{block}

    \begin{callout}{Practical note}
        The repository intentionally sticks to familiar \LaTeX{} conventions, which makes it easy to extend without learning a large custom macro layer.
    \end{callout}

\end{frame}

\section{Closing}

\makesection

\subsection{Summary}

\begin{frame}{\insertsection}{\insertsubsection}

    \begin{block}{What this template optimizes for}
        \begin{enumerate}
            \item A clean visual style that works for talks, defenses, and course presentations.
            \item A small customization surface for fast rebranding.
            \item A reusable structure that remains close to standard Beamer authoring.
        \end{enumerate}
    \end{block}

    \begin{block}{Acknowledgement}
        This project builds on ideas from the Elegant Slides template\footnote{\url{https://www.overleaf.com/latex/templates/elegant-slides/yfqyhpprvdmg}} and adapts them into a more minimal academic workflow.
    \end{block}

\end{frame}

\begin{frame}{References}
    \printbibliography
\end{frame}

\end{document}
